\documentclass[12pt,a4paper]{article}
\usepackage[utf8]{inputenc}
\usepackage[french]{babel}
\usepackage[T1]{fontenc}
\usepackage{geometry}
\geometry{top=2.5cm, bottom=2.5cm, left=2.5cm, right=2.5cm}
\usepackage{amsmath, amssymb, amsfonts}
\usepackage{graphicx}
\usepackage{booktabs}
\usepackage{xcolor}
\usepackage{hyperref}
\usepackage{listings}
\usepackage{fancyhdr}

\hypersetup{
    colorlinks=true,
    linkcolor=blue,
    filecolor=magenta,      
    urlcolor=cyan,
    pdftitle={Documentation Technique DAXX EU - OilChem Predictor},
    pdfpagemode=FullScreen,
}

\definecolor{cybercyan}{RGB}{6, 182, 212}
\definecolor{cyberindigo}{RGB}{99, 102, 241}
\definecolor{darkslate}{RGB}{30, 41, 59}

\pagestyle{fancy}
\fancyhf{}
\rhead{Projet Prédiction Chimique -- DAXX EU}
\lhead{Rapport de Conception Technique}
\rfoot{Page \thepage}

\title{\textbf{Rapport de Conception Technique \& Modélisation \\ Prédiction des Prix de Produits Chimiques\\ Plateforme Décisionnelle -- DAXX EU}}
\author{\textbf{Département Modélisation \& Data Science -- DAXX EU (Espagne)}}
\date{\today}

\begin{document}

\maketitle
\tableofcontents
\newpage

\section{Contexte Stratégique et Objectifs du Projet}
Ce projet a été conçu pour répondre aux besoins stratégiques de la direction des achats du groupe \textbf{DAXX EU} en Espagne. Dans l'industrie pétrochimique, les prix des produits finis sont extrêmement volatils et dépendent directement de l'évolution des cours de leurs composants élémentaires.

L'objectif de cette plateforme est de \textbf{prédire les prix de produits chimiques finis (cibles)} à partir de l'évolution des prix de leurs \textbf{matières premières (feedstocks)} et de leurs précurseurs en amont. En utilisant des données de marché issues d'**OilChem China**, le système identifie les dynamiques de prix, quantifie les délais de transmission des coûts au travers de la chaîne de valeur, et génère des signaux d'achats optimisés. Cet outil permet à DAXX EU d'anticiper les retournements de tendance, de fiabiliser ses prévisions budgétaires à 14 jours, et d'optimiser ses marges opérationnelles.

\section{Étapes Détaillées du Pipeline de De A à Z}
Cette section détaille chaque étape technique de la collecte de données brutes jusqu'à leur valorisation sur le dashboard.

\subsection{Étape 1 : Collecte de Données par API (oilchem\_downloader.py)}
La collecte de données s'effectue directement sur le serveur d'OilChem China via un script Python optimisé, sans émulation graphique de navigateur (\textit{headless}) :
\begin{itemize}
    \item \textbf{Identification de Session} : Le script tente de charger des cookies persistants depuis \texttt{session\_cookies.json}. Si ces cookies sont invalides ou expirés, il initie une requête POST sur l'API de connexion (\texttt{/member/login/login}) en transmettant l'identifiant et le mot de passe crypté en MD5.
    \item \textbf{Requête d'Indexation} : Pour chacun des 41 produits configurés (acétates, alcools, etc.), le script interroge l'API de recherche (\texttt{getMarketPrice}) avec des requêtes JSON contenant les caractères chinois correspondants (ex. 正丁醇 pour le n-butanol) afin de lister tous les marchés physiques disponibles.
    \item \textbf{Boucle temporelle} : Il télécharge les prix quotidiens historiques depuis le 1er janvier 2021 à l'aide de requêtes paginées vers l'API d'historique de prix.
\end{itemize}

\subsection{Étape 2 : Stockage et Consolidation des Fichiers}
Le stockage des données respecte une hiérarchie stricte :
\begin{itemize}
    \item \textbf{Données Brutes Unitaires} : Chaque série téléchargée par le downloader est stockée sous forme de fichier CSV individuel dans le dossier \texttt{data/} avec un nommage explicite contenant le nom du produit (ex. \texttt{n-Butanol\_Domestic.csv}).
    \item \textbf{Fichier Master} : Une fois le téléchargement de toutes les séries terminé, le downloader lit l'ensemble des fichiers CSV du dossier \texttt{data/}, les concatène en un seul bloc de données, et l'exporte dans le fichier master global de référence : \texttt{oilchem\_all\_data.csv}.
\end{itemize}

\subsection{Étape 3 : Nettoyage et Alignement (data\_preprocessing.py)}
Le nettoyage des données est une phase fondamentale de la finance quantitative, réalisée par le script \texttt{data\_preprocessing.py} :
\begin{itemize}
    \item \textbf{Filtrage} : Le script extrait uniquement les prix intérieurs et raffineries (codes d'activités 2 et 3) de la colonne \texttt{business\_type}.
    \item \textbf{Standardisation} : Il convertit les dates au format de référence \texttt{YYYY-MM-DD} et élimine toutes les lignes dont le prix moyen (\texttt{主流价}) est vide ou non numérique.
    \item \textbf{Séries et Pivotement} : Il génère un identifiant unique pour chaque série (\texttt{produit\_marché\_région}) et pivote la table pour passer d'un format long à un format large (les colonnes deviennent les séries de prix et les lignes les dates).
    \item \textbf{Re-Indexation et Interpolation (Forward-Fill)} : Le calendrier du marché de l'énergie comporte des interruptions (week-ends et jours fériés chinois). Pour obtenir une série quotidienne continue nécessaire aux modèles, le dataframe est réindexé sur un calendrier complet. Les valeurs manquantes sont complétées par \textit{Forward Fill} (propagation du dernier prix connu). Cette méthode est rigoureuse car elle évite le biais d'anticipation.
    \item \textbf{Export Clean} : Le dataframe aligné est exporté dans \texttt{oilchem\_aligned\_prices.csv} et \texttt{oilchem\_aligned\_prices.xlsx}.
\end{itemize}

\subsection{Étape 4 : Analyse Temporelle Lead-Lag (lead\_lag\_analysis.py)}
Le script évalue la relation temporelle de cause à effet entre les séries :
\begin{itemize}
    \item Pour chaque cible et ses intrants, il applique un décalage (lag) de 0 à 20 jours à la matière première.
    \item Il calcule le coefficient de corrélation linéaire de Pearson pour chaque lag.
    \item Le décalage optimal maximisant la corrélation absolue est exporté dans \texttt{oilchem\_lead\_lag\_results.csv}.
\end{itemize}

\subsection{Étape 5 : Moteur de Prévision et Indicateurs (financial\_prediction.py)}
\begin{itemize}
    \item \textbf{Prédiction Ridge} : Il entraîne 14 modèles de régression Ridge distincts (un par jour d'horizon) en utilisant comme variables explicatives les prix des matières premières décalés de leur lag optimal respectif.
    \item \textbf{Indicateurs Techniques} : Il calcule le Spread de marge historique, puis calcule le RSI sur 14 jours et les Bandes de Bollinger sur 20 jours de ce Spread.
    \item \textbf{Risques financiers} : Il calcule la volatilité annualisée (sur 60 jours glissants) et la Value-at-Risk (VaR 95\% sur 10 jours) du produit cible.
    \item \textbf{Saisonnalité} : Il agrège les prix moyens mensuels historiques.
    \item \textbf{JSON de Synthèse} : Toutes les informations calculées sont structurées et exportées dans \texttt{oilchem\_financial\_forecasts.json}.
\end{itemize}

\subsection{Étape 6 : Moteur de Validation (model\_backtester.py)}
Pour valider scientifiquement la pertinence économique du modèle de prédiction avant sa mise en production, un moteur de backtesting (\texttt{model\_backtester.py}) a été développé :
\begin{itemize}
    \item \textbf{Walk-Forward Backtesting} : Il simule les performances du modèle sur les données historiques en réentraînant périodiquement le modèle et en projetant les prévisions à 14 jours.
    \item \textbf{Calcul de la Précision} : Il compare la direction prédite à la direction réelle et en extrait la précision directionnelle du modèle (ex. 95.16\% pour l'Éthyl Acétate).
    \item \textbf{Calcul de l'Erreur de Prédiction} : Il calcule l'erreur absolue moyenne (MAE) en Yuans/tonne entre la prévision et le prix réel.
    \item \textbf{Simulateur d'Économies d'Achat} : Il compare la performance financière d'une stratégie d'achat basée sur les signaux du modèle (acheter quand le signal est BUY, reporter les achats quand le signal est DELAY) face à une stratégie d'achat de base régulière. Les économies réalisées (exprimées en \% et en ¥/t) sont exportées dans \texttt{oilchem\_backtest\_results.json} et affichées sur l'interface.
\end{itemize}

\subsection{Étape 7 : Automatisation et Intégration Continue (GitHub Workflows)}
Pour assurer la mise à jour quotidienne du dashboard décisionnel sans intervention humaine, le projet intègre deux workflows d'automatisation GitHub Actions :
\begin{itemize}
    \item \textbf{Mise à jour automatique (daily\_update.yml)} : Ce workflow s'exécute tous les matins à 03:00 UTC (05:00 heure espagnole). Il configure un environnement Python, charge le token secret d'OilChem depuis les secrets de sécurité de GitHub pour générer le fichier \texttt{session\_cookies.json}, exécute successivement le downloader, le preprocessing et l'analyse Lead-Lag, puis valide et pousse les modifications de données directement sur la branche principale du dépôt.
    \item \textbf{Déploiement en ligne (deploy.yml)} : Détecte chaque nouveau push sur la branche principale (y compris les mises à jour de prix matinales du bot GitHub) et compile puis déploie automatiquement l'interface utilisateur sur GitHub Pages.
\end{itemize}

\subsection{Étape 8 : Interface Utilisateur Interactive (app.js)}
Le frontend charge les fichiers de données et met en forme l'application :
\begin{itemize}
    \item \textbf{Dashboard Graphique} : Trace les courbes historiques via ApexCharts. L'utilisateur peut comparer des produits et des régions.
    \item \textbf{Visualisation de la Chaîne} : Affiche l'arborescence chimique (\textit{Chemical Value Chain}) sous le graphique principal, indiquant de manière dynamique quels intrants sont comparés.
    \item \textbf{Système d'onglets} : Structure la barre latérale en trois volets (Signaux, Achats, Matières) pour équilibrer la mise en page et éviter la surcharge visuelle.
    \item \textbf{Outils d'analyses} : Intègre un simulateur d'impact (\textit{What-If}) sur les curseurs de prix et un calculateur de budget de commandes basé sur la VaR.
\end{itemize}

\section{Modélisation Mathématique et Financière}

\subsection{Spread de Marge Théorique}
Le \textit{Spread} de transformation $S_t$ se calcule en soustrayant le coût d'achat pondéré des précurseurs du prix du produit fini :
\begin{equation}
S_t = P^{Cible}_t - \sum_{i} \left( \beta_i \times P^{Mati\grave{e}re\_i}_t \right)
\end{equation}
où $\beta_i$ représente les coefficients techniques de consommation (ex. 0.65 pour le butanol dans le Butyl Acétate).

\subsection{Bandes de Bollinger et RSI sur le Spread}
\begin{enumerate}
    \item \textbf{Bandes de Bollinger (20 jours)} :
    \begin{align}
    \mu_t &= \frac{1}{20} \sum_{k=0}^{19} S_{t-k} \\
    \sigma_t &= \sqrt{\frac{1}{20} \sum_{k=0}^{19} (S_{t-k} - \mu_t)^2} \\
    BandeSup_t &= \mu_t + 2\sigma_t \\
    BandeInf_t &= \mu_t - 2\sigma_t
    \end{align}
    \item \textbf{RSI sur 14 jours} appliqué à $S_t$ :
    \begin{equation}
    RSI_t = 100 - \frac{100}{1 + RS_t} \quad \text{avec} \quad RS_t = \frac{\text{Moyenne des hausses}}{\text{Moyenne des baisses}}
    \end{equation}
\end{enumerate}
Un signal de commande \textbf{BUY} est généré si $S_t < BandeInf_t$ ou $RSI_t < 35$ (marges comprimées, rebond attendu). Un signal \textbf{DELAY} est généré si $S_t > BandeSup_t$ ou $RSI_t > 68$ (marché surévalué).

\subsection{Décalage Lead-Lag et Pearson}
Le décalage optimal $k_{opt}$ maximisant la corrélation absolue de Pearson est donné par :
\begin{equation}
k_{opt} = \arg\max_{k \in [0, 20]} \left| \frac{\sum_{t} (X_{t-k} - \bar{X})(Y_t - \bar{Y})}{\sqrt{\sum_t (X_{t-k} - \bar{X})^2 \sum_t (Y_t - \bar{Y})^2}} \right|
\end{equation}

\subsection{Prévision Régularisée par Régression Ridge}
La régression Ridge minimise l'erreur quadratique avec une pénalité L2 sur les poids pour chaque horizon $h \in [1, 14]$ :
\begin{equation}
\min_{w_h} \| X w_h - y_{t+h} \|_2^2 + \alpha \| w_h \|_2^2
\end{equation}
La contrainte L2 stabilise le vecteur de poids $w_h$, ce qui empêche le modèle de sur-réagir à la multicollinéarité inhérente aux prix de la même chaîne chimique.

\subsection{Calcul de Risque et Volatilité}
\begin{itemize}
    \item \textbf{Volatilité Logarithmique Annualisée} :
    \begin{equation}
    \sigma_{annuelle} = \text{std}(\ln(P_t / P_{t-1})) \times \sqrt{252} \times 100
    \end{equation}
    \item \textbf{Value-at-Risk (VaR 95\% 10 jours)} :
    \begin{equation}
    VaR_t = 95^{\text{ème}}\text{ centile de } (P_{\tau} - P_{\tau-10})
    \end{equation}
\end{itemize}

\section{Difficultés Techniques Réelles Résolues}
Les véritables verrous technologiques et scientifiques qui ont été levés lors de ce projet sont les suivants :

\begin{enumerate}
    \item \textbf{Le Verrou de l'Authentification de l'API (Format Custom du Token Header)} :
    Lors de la connexion automatisée à la base de données d'OilChem China, les requêtes étaient systématiquement rejetées par le serveur malgré des identifiants et des cookies valides. Après analyse réseau et rétro-ingénierie, nous avons découvert que le serveur exigeait un en-tête HTTP personnalisé nommé \texttt{token} contenant la valeur brute sous la forme \texttt{\_member\_user\_tonken\_=<valeur\_du\_token>}. L'implémentation de cette syntaxe a permis d'extraire les données directement via des appels d'API ultra-rapides, évitant l'utilisation d'une infrastructure lourde de navigateur virtuel (Selenium).
    
    \item \textbf{La Multicollinéarité des Prix en Cascade} :
    Les prix d'une même chaîne chimique étant intrinsèquement corrélés (le brut influençant le naphta, qui influence la matière première, qui influence la cible), une régression linéaire classique (\textit{OLS}) sur-apprenait et produisait des coefficients aberrants. Le passage à des modèles de **Régression Ridge** avec régularisation L2 a stabilisé les coefficients de sensibilité, garantissant la fiabilité des prévisions de prix à 14 jours.
    
    \item \textbf{L'Instabilité Aléatoire de l'Ordre des Régions} :
    Lors des phases de mise à jour du modèle, l'ordre des régions disponibles à l'ouverture de l'application changeait de façon imprévisible (passant de *East China* à *South China* par exemple) sans modification du code de l'interface. L'origine de cette anomalie résidait dans l'utilisation de structures Python \texttt{set} pour répertorier les régions dans le backend. En Python, les sets sont non ordonnés et dépendent du hachage mémoire, aléatoire d'une exécution à l'autre. Le problème a été résolu en appliquant un tri alphabétique strict à la liste des régions avant son exportation dans le fichier de configuration JSON.
    
    \item \textbf{La Résistance au Cache du Navigateur Client} :
    L'application d'améliorations visuelles (comme la réorganisation de la barre latérale en onglets) ne se reflétait pas sur les postes clients en raison d'une mise en cache agressive des fichiers CSS et JS par les navigateurs. Ce problème récurrent a été contourné en implémentant un mécanisme de cache-busting basé sur des variables de requêtes explicites (\texttt{?v=42}) incrémentées à chaque déploiement.
\end{enumerate}

\section{Limites du Modèle Prédictif en Période de Crise}
Bien que le modèle Ridge soit performant avec une précision directionnelle historique élevée, il présente des limites intrinsèques en période d'anomalie de marché (chocs géopolitiques majeurs ou ruptures d'approvisionnement brutales) :
\begin{itemize}
    \item \textbf{Absence de mémoire des extrêmes} : Le modèle s'entraîne sur des corrélations moyennes. Lors d'une hausse parabolique anormale suivie d'une prédiction immédiate, le modèle peut générer des signaux d'achat prématurés en pensant que le bas de cycle est atteint, alors que la dynamique baissière du marché se poursuit.
    \item \textbf{Recommandation} : Il est recommandé de coupler la décision algorithmique avec une observation humaine et de ne valider les achats que lorsque l'indicateur de direction à 14 jours (\texttt{14D Trend}) confirme une stabilisation ou un retournement haussier réel.
\end{itemize}

\end{document}
